\documentclass[11pt,a4paper]{article}
\newcommand{\ExperimentNumber}{036}
\newcommand{\AppendixFooter}{Fixed EVA factors $\cdot$ G-set benchmark}
\usepackage{maxcut-appendix}
\newcommand{\PrincipalFinding}{The centroid sweep reached a mean gain of 4.384\%, versus 4.239\% for random sweep (+0.145 pp). All three cut-guided variants improved on random sweep in the equal-graph mean; inertia-only search did not.}
\newcommand{\ConvergenceFinding}{Repeated GW achieves 2.411\%, random sweep 4.239\%, and the centroid, antipodal and hybrid variants 4.384\%, 4.383\% and 4.373\%, respectively. Gains are measured relative to the common initial cut.}
\newcommand{\EffectFinding}{Centroid: +0.145 pp, interval [+0.086, +0.209]; antipodal: +0.144 pp, interval [+0.095, +0.197]; hybrid: +0.134 pp, interval [+0.087, +0.184]. These intervals exclude zero, but do not establish a ranking among the three guided variants.}
\newcommand{\CostFinding}{Mean archived seconds are 0.163 for random sweep and 0.197 for centroid sweep; inertia sweep takes 1.069. The latter yields lower cut quality despite examining a similar number of intervals.}
\newcommand{\RecordFinding}{Random sweep registers 17.31 strict records on average, compared with 16.13 for centroid sweep. Spherical two-means stops after 24.01 attempts on average. More records do not imply a better final incumbent.}
\newcommand{\GraphFinding}{Centroid wins on 52 of 71 graphs, antipodal on 55, and hybrid on 50. Their respective losses (14, 13, 18) rule out universal dominance.}
\newcommand{\GeometryFinding}{Direct inertia minimization attains 2.839\% versus 4.239\% for random sweep (-1.400 pp). Spherical two-means achieves 2.813\%. The cut-guided variants perform better on average than these geometry-only updates.}
\newcommand{\AuditFinding}{The completed audit covered 716,826 records, including 487,354 new and 229,472 reused records; 48,060 sampled event intervals were archived. All methods recorded zero invalid candidate flags.}
\newcommand{\ResultRows}{
Repeated GW & 2.411 & -1.828 & [-2.151, -1.509] & 0/2/69 \\
Random sweep & 4.239 & +0.000 & [+0.000, +0.000] & 0/71/0 \\
Spherical two-means & 2.813 & -1.425 & [-1.725, -1.137] & 1/2/68 \\
\textbf{Centroid sweep} & 4.384 & +0.145 & [+0.086, +0.209] & 52/5/14 \\
Antipodal sweep & 4.383 & +0.144 & [+0.095, +0.197] & 55/3/13 \\
Inertia sweep & 2.839 & -1.400 & [-1.699, -1.116] & 1/2/68 \\
Hybrid sweep & 4.373 & +0.134 & [+0.087, +0.184] & 50/3/18 \\
}

\hypersetup{pdftitle={Experiment 036: Geometry-guided hyperplane rounding}}
\begin{document}
\appendixheading{EXPERIMENT 036 \enspace / \enspace EMPIRICAL RESULTS}{Geometry-guided hyperplane rounding}{Numerical appendix \quad | \quad G-set benchmark \quad | \quad 26 November 2026}
\textbf{Research question.} Can geometry-informed rotations improve the original signed Max-Cut objective relative to random-tangent hyperplane sweeps, and is directly minimizing inertia an effective proxy for maximizing the cut?

\textbf{Principal finding.} \PrincipalFinding

\section{Experimental design}
All seven methods start from identical Goemans--Williamson (GW) normals, using fixed EVA semidefinite factors on 71 graphs with 16 paired starts per graph. Each sweep has a budget of 100 attempts; spherical two-means may stop earlier. No new SDP solve or discrete local refinement was performed. The two baseline trajectories are reused from Experiment 033.

For a signed weighted graph $G=(V,E,w)$, let $y_i\in\mathbb R^p$ be the unit embedding of vertex $i$. A unit normal $r$ induces $x_i=\mathbf 1\{y_i^\top r>0\}$. The objective and normalized spherical inertia are
\begin{equation*}
 C(x)=\sum_{(i,j)\in E}w_{ij}\mathbf 1\{x_i\ne x_j\},\qquad
 I_{\mathrm{sph}}(x)=2-\frac{2}{n}\left(\|S_0\|+\|S_1\|\right),
\end{equation*}
where $S_k=\sum_{i:x_i=k}y_i$. A rotation $r(\theta)=r\cos\theta+v\sin\theta$ with $v^\top r=0$ induces a finite sequence of partitions. The random, centroid, antipodal and hybrid cut sweeps optimize the signed cut along the proposed arc. The geometry-only arms update by spherical inertia and separately retain the best cut encountered.

\section{Methods compared}\label{sec:methods}
A hyperplane normal determines a partition. A sweep tests the distinct partitions encountered as the normal rotates; the cut is constant between vertex-crossing events. Let $z\sim\mathcal N(0,I_p)$, let $u(a)=a/\|a\|$ for $a\ne0$, and let $P_r=I_p-rr^\top$ project onto the tangent space at $r$. The two geometry-derived targets are
\[
 q_c=u\bigl(u(S_1)-u(S_0)\bigr),\qquad q_a=u(S_1-S_0).
\]
The first gives equal weight to the two normalized cluster centers; the second uses their unnormalized sums. The seven methods differ in how they choose a candidate and decide whether to move.

\begin{description}
\setlength{\itemsep}{2pt}\setlength{\parsep}{0pt}
\item[Repeated GW.] Draw a fresh independent normal $r'=u(z)$ at each attempt and evaluate its partition. Keep the highest cut seen; subsequent draws do not depend on it. This is the standard repeated-random-rounding baseline.
\item[Random sweep.] At the current best normal $r$, choose the random tangent $v=u(P_r z)$. Enumerate the partitions along the semicircle $r(\theta)$ and select the largest signed cut, retaining the incumbent if no improvement is found.
\item[Spherical two-means.] Normalize each cluster sum to form two unit centers; replace $r$ directly by $q_c$, which assigns vertices to the closer unit center. Accept a strict reduction of spherical inertia and stop on the first non-improving attempt. No rotation sweep is performed.
\item[Centroid sweep.] Use $v=u(P_r q_c)$ to guide the rotation toward the difference of the two normalized centers, then choose the largest signed cut along that arc. Geometry chooses the search direction; the graph objective chooses the update.
\item[Antipodal sweep.] Use $v=u(P_r q_a)$ and maximize the signed cut along the arc. The target $q_a$ is the optimal direction for fixed labels if the two unit centers are constrained to be antipodal. Unlike $q_c$, it weights cluster contributions by their resultant lengths.
\item[Inertia sweep.] Propose the centroid tangent $u(P_r q_c)$, but select the interval that minimizes spherical inertia, rather than maximizing the cut. Accept only strict inertia descent. This tests whether geometric compactness alone is a useful surrogate for the original objective.
\item[Hybrid sweep.] Cycle through centroid, antipodal and random tangent proposals. At each attempt maximize the signed cut along the arc and accept only a strict cut improvement. The mixture combines geometric guidance with unguided exploration.
\end{description}

\textbf{Shared search rules.} Sweep proposals use $-\pi/2<\theta<\pi/2$ and keep the incumbent as a fallback. In the new guided sweeps, an unavailable or degenerate guided tangent triggers a random tangent; after a rejected update, the next proposal is random. The cut-guided methods advance on strict cut improvement, while spherical two-means and inertia sweep advance on strict inertia reduction. All methods report the best cut encountered, even if a geometry-only candidate is not accepted into the search state.


\section{Primary numerical results}
\begin{table}[H]\centering\small
\caption{Equal-graph mean gain and paired difference versus random sweep. The interval is a 95\% graph-bootstrap interval; W/T/L counts graphs won, tied and lost.}\label{tab:results}
\begin{tabular*}{\linewidth}{@{\extracolsep{\fill}}lrrrl@{}}
\toprule
Method & Gain (\%) & $\Delta$ (pp) & 95\% interval (pp) & W/T/L\\
\midrule
\ResultRows
\bottomrule
\end{tabular*}
\end{table}
\smallskip
{\small Gains are changes from the initial GW cut divided by the graph's mean initial GW cut. Each graph receives equal weight, and intervals use 10,000 graph-level bootstrap resamples; they are not multiplicity-adjusted.}

\sectionpage{Convergence and paired effects}
\resultfigure[.30\textheight]{convergence}{Mean incumbent gain over 100 attempts. Stopped spherical two-means trajectories are held constant.}{\ConvergenceFinding}
\resultfigure[.32\textheight]{paired_effects}{Paired graph-bootstrap effects relative to random sweep. The vertical reference is zero; horizontal bars are 95\% intervals.}{\EffectFinding}

\sectionpage{Computational effort and record frequency}
\resultfigure[.30\textheight]{quality_cost}{Mean gain against archived method seconds (A) and examined intervals (B). Axes use a symmetric logarithmic scale.}{\CostFinding}
\resultfigure[.29\textheight]{records}{Mean number of strict cut records (A) and attempted updates (B). A cut record is an increase in the best signed cut, not a reduction in inertia.}{\RecordFinding}

\sectionpage{Graph heterogeneity and geometry}
\resultfigure[.32\textheight]{graph_heatmap}{Graphwise advantage relative to random sweep, for each new method. Positive values favor the new method.}{\GraphFinding}
\resultfigure[.29\textheight]{inertia_trajectories}{Active-candidate spherical (A) and Euclidean (B) inertia, normalized by the number of vertices.}{Late active-only means for two-means are selected by survival. A decrease in inertia is not a guarantee of a larger signed cut.}

\sectionpage{Inertia and the original cut objective}
\resultfigure[.49\textheight]{inertia_cut_density}{Candidate gain versus spherical-inertia reduction relative to the common starting geometry; logarithmic hexbin density. Panels correspond to the five new methods.}{\GeometryFinding}
\textbf{Implication.} Geometric information is useful as a proposal mechanism, but the acceptance objective matters. Pooled candidate-level densities are exploratory, because starts and iterations within a graph are dependent and graphs differ in scale.

\sectionpage{Search diagnostics}
\resultfigure[.53\textheight]{search_diagnostics}{Strict record probability, Hamming turnover, projective angular displacement and minority-side fraction, by active attempt. Global sign reversal is ignored in turnover and angle.}{Record rates and movement characterize the trajectory rather than establish superiority. Candidate curves may differ from incumbent curves; an inertia-descent arm can reject a candidate yet retain its cut independently.}
\textbf{Interpretive caution.} Timing comparisons are descriptive because the archived baselines and new methods came from different runs. The recorded time excludes some diagnostics and is not an equal-time benchmark.

\sectionpage{Rotation-event landscapes}
The rotation profiles below show the first arc at paired seed zero. Each point corresponds to an open interval between hyperplane crossing events. Blue denotes signed cut; orange denotes spherical inertia. The plotted profiles are sampled diagnostics, not independent observations.
\resultfigure[.37\textheight]{event_G1}{Event landscape for G1, four guided rotation variants.}{}
\resultfigure[.37\textheight]{event_G11}{Event landscape for G11, four guided rotation variants.}{}
\clearpage
\resultfigure[.47\textheight]{event_G14}{Event landscape for G14, four guided rotation variants.}{}
\textbf{Interpretation.} The extrema of the two objectives need not coincide. Choosing the arc direction from the geometry and selecting the partition by its original signed cut are distinct operations. These three examples illustrate that distinction; they do not establish a population-level association.
\section{Verification and reproducibility}
\AuditFinding\ The independent audit reconstructs partitions from saved normals, evaluates the signed objective using the original graph edges, and checks incumbent monotonicity, geometry, checksums and factor identities. The centroid-angle comparison uses cosine space to avoid endpoint ill-conditioning. All results in this appendix were generated from the existing audited archive, without repeating any optimization. The source project includes the original figure PDFs and compact summary data; the full trajectories are retained separately under the VM archival policy.

\sectionpage{Graphwise convergence}
The following six figures cover all 71 graphs in benchmark order. Each curve averages the 16 paired starts; vertical scales are graph-specific. Stopped two-means trajectories are held constant. Compare methods within a panel rather than absolute heights across panels.
\resultfigure[.77\textheight]{graph_curves_1}{Graphwise convergence, batch 1 of 6.}{}
\clearpage
\resultfigure[.86\textheight]{graph_curves_2}{Graphwise convergence, batch 2 of 6.}{}
\clearpage
\resultfigure[.86\textheight]{graph_curves_3}{Graphwise convergence, batch 3 of 6.}{}
\clearpage
\resultfigure[.86\textheight]{graph_curves_4}{Graphwise convergence, batch 4 of 6.}{}
\clearpage
\resultfigure[.86\textheight]{graph_curves_5}{Graphwise convergence, batch 5 of 6.}{}
\clearpage
\resultfigure[.86\textheight]{graph_curves_6}{Graphwise convergence, batch 6 of 6.}{}
\end{document}
